\section{Conclusion}\label{sec:conclusion}

This paper presented a parameter-efficient five-stage convolutional
architecture for binary and multi-class oral disease classification. The
architecture departs from a standard EfficientNetV2-B0 in two deliberately
minimal ways: a custom AttentionHub replaces the donor Block-4, and the
redundant tail (Block-6 together with the convolutional head) is removed.
Through a seven-cell ablation of the AttentionHub, reinforced by a confirmatory
negative-result run, we derived a \textbf{role-complementarity principle} that
constrains how attention modules should be combined in this setting: pair
modules with disjoint functional roles, or pay a measurable accuracy penalty.
The principle forced a sequential Triplet\arrowto{}SE cascade as the proposed
\emph{AttentionHub-v2}, which under a matched fair training recipe attains
$99.06\,\%$ binary accuracy and $99.51\,\%$ subtype accuracy on a seven-class
held-out test set, the highest subtype score in a benchmark of nine standard
CNN and Transformer baselines, at $4.79$~M parameters and $0.493$~GFLOPs,
making it $2.1\times$ smaller than EfficientNet~V2-B2 and $5.0\times$ smaller
than Inception~V3, the two strongest baselines. Grad-CAM++ and LIME panels
indicate that the model attends to lesion tissue rather than to incidental
visual cues; this interpretability evidence is qualitative and is consistent
with, rather than a formal verification of, the model's clinical alignment. The
full benchmark, the ablation grid, and the explainability artifacts are released
to support reproducibility and to make the role-complementarity principle
available for re-testing on adjacent medical-imaging distributions.
